% Chapter II, §18: Experiments on the onset of thermal instability in fluids
% PDF pages 76--78, book pages 59--61
% Covers: §18 introduction + (a) Bénard's experiments + start of (b) Schmidt–Milverton principle
% Continues from chapter_2_part2.tex (equations up to 322)

%%---------------------------------------------------------------------
%% §18.  Experiments on the onset of thermal instability in fluids
%%---------------------------------------------------------------------

\section{Experiments on the onset of thermal instability in fluids}\label{sec:18}

In this section we shall give an account of some of the experimental
work on the onset of thermal instability in fluids which bears on the
theoretical developments.

As we stated at the outset, and as the subtitle to this chapter indicates,
the experiments which stimulated this field of inquiry are those of
B\'enard.  Indeed, it was to account for the `interesting results obtained
by B\'enard's careful and skilful experiments' that Rayleigh undertook
his theoretical investigation; the quotation is in fact the opening remark
in Rayleigh's paper.

%%---------------------------------------------------------------------
%% (a) Bénard's experiments
%%---------------------------------------------------------------------

\subsection*{(a) B\'enard's experiments}

B\'enard carried out his experiments on very thin layers of fluid, about
a millimetre in depth, or less, standing on a levelled metallic plate

\figurePlaceholder{10}{The original apparatus of B\'enard.}

\noindent
maintained at a constant temperature (see Fig.~10).  The upper surface
was usually free and being in contact with the air was at a lower temperature.
B\'enard experimented with several liquids of differing physical
constants.  He was particularly interested in the r\^ole of viscosity; and
as liquids of high viscosity he used melted spermaceti and paraffin.  In
all cases, B\'enard found that when the temperature of the lower surface
was gradually increased, at a certain instant, the layer became reticulated
and revealed its dissection into cells.  He further noticed that there were
motions inside the cells: of ascension at the centre, and of descension
at the boundaries with the adjoining cells.  B\'enard distinguished two
phases in the succeeding development of the cellular pattern: an initial
phase of short duration in which the cells acquire a moderate degree
of regularity and become convex polygons with four to seven sides and
vertical walls; and a second phase of relative permanence in which the
cells all become equal, hexagonal, and properly aligned.  Fig.~1, which is
a reproduction of one of B\'enard's original photographs, illustrates the
remarkable regularity which can be attained.  B\'enard himself was so
interested in observing the emergence of the second phase and the
associated motions that he was not very explicit regarding the conditions
necessary for the onset of instability.  However, in a later analysis of his
early observations, he states that they are in qualitative agreement with
the requirements of Rayleigh's criterion.

Some doubts have recently been raised whether, in view of the thinness
of the layers of fluid with which B\'enard experimented, his observations
can all be accounted for in terms of effects arising from surface
tension.  As B\'enard himself pointed out quite explicitly in his early papers,
effects of surface tension are undoubtedly present; but they do not
appear to be primary in most of his experiments.

%%---------------------------------------------------------------------
%% (b) The Schmidt–Milverton principle (beginning)
%%---------------------------------------------------------------------

\subsection*{(b) The Schmidt--Milverton principle for detecting the onset of thermal instability}

Of the vast literature on experiments relating to thermal convection,
we shall refer only to that which bears directly on the detection of the
onset of thermal instability, and the determination of the critical
Rayleigh number.  The earliest of the experiments in this latter category
appear to be those of Schmidt and Milverton.  In their experiments,
Schmidt and Milverton incorporated a principle for the detection of the
onset of thermal instability which is so direct and simple that it has
served as the basis for all later experiments in this subject.  In view of
the importance of this principle, we shall explain its origin and the
manner of its use.

Suppose that a difference of temperature $\Delta T$ $(= T_1 - T_2;\; T_2 > T_1)$
is maintained between the two sides of a layer of fluid of depth~$d$.  In a
steady state, this implies that there is a constant outward flux of heat
from the surface at the higher temperature.  Let $\mathscr{Q}$ denote this flux
of heat: it measures the quantity of heat (in calories) emerging from the
surface, per unit area and per unit time.  If the flow of heat across the
fluid layer is entirely by conduction, then clearly
\begin{equation}
\label{eq:2-323}
\mathscr{Q} = k\,\Delta T / d,
\end{equation}
where $k$ is the coefficient of heat conduction.
